\chapter{Evaluation Results}

Before arriving at gradient accumulation with percentile-based sparsification, we explored two alternative compression strategies based on summary statistics and information-theoretic filtering. Neither approach achieved competitive performance; details are provided in Appendix~\ref{app:preliminary-compression}.

\section{Gradient Accumulation Results}

All experiments were conducted on a machine equipped with an NVIDIA RTX 3090 GPU (24~GB VRAM) running Ubuntu Linux. We evaluate our gradient accumulation approach across three configurations on CarRacing-v3: (1) baseline pipeline parallelism without compression, (2) 90th percentile sparsification without decay ($\lambda = 1.0$), and (3) 90th percentile sparsification with 0.99 buffer decay. All runs use the proven hyperparameter configuration: \texttt{num\_envs=1}, \texttt{num\_steps=4096}, learning rate $4 \times 10^{-5}$, 32 minibatches, and a total training budget of 1,000,000 timesteps (244 iterations). Gradient compression activates after a 30,000-step warm start period. Episodic returns are smoothed using exponential moving average (EMA) with decay 0.9.

\section{Training Performance}

Figure~\ref{fig:exp4-returns} shows training performance across the three configurations. The vertical dashed line marks the warm start boundary at 30,000 steps, after which gradient compression activates.

During the warm start phase (0--30,000 steps), all three runs follow identical trajectories, as expected since no compression is applied. After compression activates, the configurations begin to diverge.

The baseline without compression reaches an EMA return of approximately 496 before crashing at iteration 169. The 90th percentile configuration without decay follows a similar early trajectory but shows slower improvement after 200,000 steps, reaching an EMA return of 304 at the same cutoff. The 90th percentile with 0.99 decay achieves the strongest performance, climbing steadily to an EMA return of 719, surpassing the baseline by 45\%.

\begin{figure}[htbp]
\centerline{\includegraphics[width=\textwidth]{images/exp4_episodic_returns.pdf}}
\caption{Training performance comparison over approximately 692,000 steps. The 90th percentile with 0.99 decay outperforms both the baseline and the no-decay variant. The dashed vertical line indicates warm start completion at 30,000 steps.}
\label{fig:exp4-returns}
\end{figure}

\section{Policy Evaluation}

\begin{figure}[htbp]
\centerline{\includegraphics[width=\textwidth]{images/exp4_performance_comparison.pdf}}
\caption{Evaluation return distributions over 100 independent episodes per checkpoint. Box plots show median, IQR, and whiskers (1.5$\times$IQR). The decay variant matches baseline performance, while the no-decay 90th percentile achieves only 53\% of baseline mean return.}
\label{fig:exp4-boxplot}
\end{figure}

To characterize final policy quality independent of training curve smoothing, we evaluated the last saved checkpoint from each configuration over 100 independent episodes of CarRacing-v3 with a fixed random seed. Each checkpoint was loaded and rolled out without any further training updates. The baseline was evaluated at iteration 160 (the last checkpoint before its crash at iteration 169), the 90th percentile at iteration 170 (last before crash at 172), and the decay variant at iteration 180 (last before manual stop at 121).

Figure~\ref{fig:exp4-boxplot} shows box plots of episodic returns for the three configurations. The baseline achieves a mean return of 534.63 ($\pm$196.78) with a median of 534.73. The 90th percentile without decay achieves a mean of 285.56 ($\pm$179.78) with a median of 266.95, approximately 53\% of baseline performance. The 90th percentile with 0.99 decay achieves a mean of 572.52 ($\pm$226.58) with a median of 584.48, slightly exceeding baseline performance at 107\% of the baseline mean.

The interquartile ranges tell a similar story: the baseline spans [357, 721], the 90th percentile spans [190, 365], and the decay variant spans [417, 786]. The decay variant's IQR is shifted above the baseline's, consistent with the training curve observations. The 90th percentile without decay shows a tighter distribution but at substantially lower returns, with two outlier episodes above 800 that suggest the policy occasionally finds good trajectories but cannot do so consistently.

These evaluation results reinforce the training performance findings: buffer decay mitigates the staleness problem in gradient accumulation, producing a policy comparable to the uncompressed baseline. Without decay, transmitting only 10\% of gradients leads to a meaningful degradation in final policy quality.

\section{Communication Overhead}

A primary goal of gradient compression is to reduce the volume of data transferred between machines during training. Figure~\ref{fig:exp4-transfer} shows cumulative data transfer over the full training run for each configuration.

The baseline transmits the full $(b \times 4096)$ gradient tensor from Machine~1 to Machine~0 every minibatch, accumulating approximately 266~GB over the training run. Both compressed configurations transmit only a sparse representation containing the top 10\% of gradients by magnitude, along with their indices and shape metadata. This reduces cumulative transfer to approximately 173~GB, a 35\% reduction.

The two compressed variants show nearly identical transfer volumes. This is expected: the decay factor $\lambda$ controls \emph{which} gradients exceed the percentile threshold, not \emph{how many}. Both configurations apply the same 90th percentile cutoff, so approximately the same fraction of gradient values is transmitted each minibatch. The difference between the two lies entirely in gradient quality: decay ensures that transmitted gradients reflect recent optimization state rather than stale accumulations.

\begin{figure}[htbp]
\centerline{\includegraphics[width=\textwidth]{images/exp4_network_transfer.pdf}}
\caption{Cumulative data transfer over the training run. Both compressed configurations reduce transfer by approximately 35\% relative to the uncompressed baseline. The dashed vertical line indicates warm start completion.}
\label{fig:exp4-transfer}
\end{figure}

Table~\ref{tab:transfer-summary} summarizes the transfer volumes and savings for each configuration alongside final policy performance.

\begin{table}[htbp]
\centering
\caption{Communication overhead and policy performance summary.}
\label{tab:transfer-summary}
\begin{tabular}{lccc}
\toprule
\textbf{Configuration} & \textbf{Transfer (GB)} & \textbf{Reduction} & \textbf{Mean Return} \\
\midrule
No Compression          & 266 & --  & 534.63 \\
90th Percentile         & 173 & 35\% & 285.56 \\
90th Percentile + Decay & 173 & 35\% & 572.52 \\
\bottomrule
\end{tabular}
\end{table}

The 90th percentile with decay achieves the same 35\% reduction in communication as the no-decay variant while producing a policy that matches or exceeds baseline performance. This demonstrates that gradient accumulation with percentile-based sparsification and buffer decay can meaningfully reduce communication overhead in pipeline-parallel reinforcement learning without sacrificing policy quality.

\section{Effect of Buffer Decay}

The decay factor $\lambda = 0.99$ discounts the accumulation buffer before adding new gradients each minibatch. With 32 minibatches per epoch and 10 update epochs, the decay is applied 320 times per training iteration, yielding a per-iteration retention factor of $0.99^{320} \approx 0.04$. This means gradients that do not exceed the percentile threshold within a single iteration are effectively discarded.

The improved performance of the decay variant over the no-decay 90th percentile can be explained by the staleness problem in the accumulation buffer. Without decay, below-threshold gradients accumulate slowly across iterations. By the time they eventually cross the percentile threshold, the optimization landscape has shifted. These stale gradients now point in a direction that is no longer optimal, causing the CNN feature extractor to take harmful update steps. The decay factor mitigates this by discounting older gradients, ensuring that only recently significant gradient signal reaches the transmission threshold. In effect, the decay acts as a relevance filter: gradients that are not important enough to send immediately are progressively forgotten rather than allowed to accumulate into a potentially misleading update.

\section{Gradient Compression Without Decay}

The 90th percentile configuration without decay ($\lambda = 1.0$) reached an EMA return of 304, which corresponds to about 61\% of the baseline performance. This result is consistent with our prior observations: sending only 10\% of the gradient values reduces the quality of the learning for the CNN feature extractor. That is because the remaining 90\% of gradient dimensions do not receive any update until each individual dimension accumulates enough to cross the threshold. As can be seen by examining the difference in performance between compressed and uncompressed training over time, it appears that the cumulative effect of errors due to missing gradient dimensions increases as training proceeds.

In addition, we ran a series of additional experiments using 0.99 decay with higher percentile thresholds. With the 95th percentile (i.e., sending only 5\% of the gradients), the agent does appear to learn; however, at a much reduced rate than when all gradients were sent. When we used the 99th percentile (i.e., sending only 1\% of gradients), the agent did not learn at all, confirming that there is a minimum amount of gradients required to provide a sufficiently strong learning signal for the CNN feature extractor to perform well. In order to ensure that the baseline and 90th percentile with decay were reproducible across different seeds, we ran both again using a different random seed. Both converged to similar levels of performance. The complete training curves for all five configurations are provided in Appendix~\ref{app:percentile-sweep}.
